\Fidel{This subsection formally introduces and defines SPCs.
\begin{enumerate}
    \item Define the multi-time Pauli twirling operation.
    \item Theorem: Twirling projects a quantum process tensor to a stochastic (correlated) Pauli error distribution.
    \item Formally define SPCs.
    \item Worked examples?
\end{enumerate}
}
% \subsection{Formally define SPCs from twirled process tensors}
\begin{definition}\label{def:multitime-twirl}
    Define multi-time Pauli twirl.
\end{definition}

\begin{lemma}\label{lem:multi-time-twirl-cptp}
    The multi-time Pauli twirl operator is a local CPTP superprocess (or a super-superchannel in the language of higher order quantum maps \cite{tarantohigherorder2025}).
\end{lemma}


% \subsection{First theorem}
\begin{theorem}\label{thm:pt-to-spc}
  Pauli twirling a multi-time process removes all quantum temporal correlations, resulting in a spatiotemporal Pauli channel.
\end{theorem}


\begin{definition}\label{def:def-spc}
    Define SPCs. (Maybe make this definition 1?)
\end{definition}

Spatiotemporal Pauli string $\mathbf{P}\in\mathbb{P}^{(n,k)}$:
$\mathbf{P} = \bigotimes_{j=0}^{k-1}P_j$. \\
A Spatiotemporal Pauli Channel (SPC) $\mathbf{\Lambda}_\mathbf{P}$ is entirely characterised by a probability distribution over the set of SPCs, $\mathbb{P}^{(n,k)}$.
\begin{equation}
    \Pr{\mathbf{P}} \geq 0, \quad \sum_{\mathbf{P}\in\mathbb{P}^{(n,k)}} \Pr{\mathbf{P}} = 1
\end{equation}
The Choi matrix representation of $\mathbf{\Lambda}$ is the convex sum of Choi matrices for each deterministic Pauli string process $\Upsilon_{\bm{\Lambda}}$ weighted by their respective probabilities
\begin{equation}
    \Upsilon_{\bm{\Lambda}} = \sum_{\mathbf{P}\in\mathbb{P}^{(n,k)}} \Pr{\mathbf{P}} \Upsilon_{\mathbf{P}}
\end{equation}
\begin{equation}
    \Upsilon_{\mathbf{P}} = \bigotimes_{j=0}^{k-1}\Pi_j
\end{equation}
$\Pi_j = \kketbbra{P_j}{P_j}$ and $
\{\kket{I}, \kket{X}, \kket{Y}, \kket{Z}\} = \{\ket{\Phi^+}, \ket{\Phi^-}, \ket{\Psi^+}, \ket{\Psi^-}\}
$

\begin{equation}\label{eq:twirl-upsilon}
    \mathcal{T}_P^{(k)}(\Upsilon_{k:0}) = \frac{1}{|\tilde{\mathbb{P}}^{(n,k)}|}\sum_{\tilde{\mathbf{P}} \in \tilde{\mathbb{P}}^{(n,k)}}\tilde{\mathbf{P}}\Upsilon_{k:0}\tilde{\mathbf{P}}^\dagger
\end{equation}
where $\tilde{\mathbf{P}} = \bigotimes_{j=0}^{k-1}(P_j^T\otimes P_j) = \bigotimes_{j=0}^{k-1} \hat{P}_j$ is the time-ordered tensor product of Pauli superoperators. $|\tilde{\mathbb{P}}^{(n,k)}|=4^{nk}$.

\begin{equation}\label{eq:spc-probabilities}
    \Pr{\mathbf{P}} = \Tr\left[\left(\bigotimes_{j=0}^{k-1}\Pi_j\right)\Upsilon_{k:0}\right]
\end{equation}

\subsection*{Validity of Pauli process tensor}
\begin{enumerate}
    % \item Probabilities: $\sum_$
    \item Positivity constraint: $\Upsilon_{k:0}\geq0$\\
    This condition is fulfilled for $\Upsilon_{k:0}^{\mathcal{T}_P}$ as seen in \cref{eq:ppt-bell-basis}.
    \item Affine constraint: $\Tr_\out[\Upsilon_{k:0}] = \mathbb{I}_\inp \otimes \Upsilon_{k-1:0} \forall k$ \\
    Denoting $\braket{P_{\mu_1}\cdots P_{\mu_{2k}}} \coloneq \Tr{[(P_{\mu_1}\otimes \cdots\otimes P_{\mu_{2k}})\Upsilon_{k:0}]}$, this condition can be rewritten as
    \begin{equation}
        \begin{split}
            &\braket{P_{\inp_1}^{\mu_1}\cdots P_{\out_{k-1}}^{\mu_{2k-2}}\tilde{P}_{\inp_k}^{\mu_{2k-1}}\mathbb{I}_{\out_k}^{\mu_{2k}}} = 0 \\
            &\braket{P_{\inp_1}^{\mu_1}\cdots \tilde{P}_{\inp_{k-1}}^{\mu_{2k-3}}\mathbb{I}_{\out_{k-1}}^{\mu_{2k-2}}\mathbb{I}_{\inp_k}^{\mu_{2k-1}}\mathbb{I}_{\out_k}^{\mu_{2k}}} = 0 \\
            &\vdots \\
            &\braket{\tilde{P}_{\inp_1}^{\mu_1}\mathbb{I}_{\out_1}^{\mu_{2}}\cdots \mathbb{I}_{\inp_k}^{\mu_{2k-1}}\mathbb{I}_{\out_k}^{\mu_{2k}}} = 0,
        \end{split}
    \end{equation}
where $\tilde{P}\in\{X,Y,Z\}$. This condition is naturally fulfilled \cref{eq:ppt-pauli-basis}.

\textcolor{red}{Angus: these affine constraints might make sampling more efficient -- you don't need the full process tensor for times $t > k$ if you only want samples of the first $k$ timesteps, because all the partial traces over the latter timesteps simply give identities.}
\end{enumerate}


% \subsection{Mathematical properties}
% \begin{enumerate}
%     \item Are Pauli process tensors equivalent to pseudo-density matrices up to some conjugation?
%     \item Pauli process tensor $\chi$-matrix positivity conditions?
%     \item Can the Pauli process tensor be expressed by a classical distribution?
%     \item When does a process tensor map to a Markovian distribution after twirling?
%     \item Can we bound the correlations of process tensor from the twirled process 
%     tensor?
% \end{enumerate}
% \subsection{Interpretation}
% \subsection{Pauli process tensor metrics}
% \begin{enumerate}
%     \item Pairwise correlation as a function of spacetime distance
%     \item Based on the underlying process (size of environment, locality of interaction Hamiltonian, quantum correlations) what can we say about the space-time correlations of Pauli channels?
    
% \end{enumerate}
% \subsection{Tomographic construction?}